\section*{Capítulo \ref{ch:b1:logic} --- Lógica, Conjuntos e Aplicações}

\begin{solution}{exo:b1:logic:1}
Negações, empurrando $\lnot$ através de cada quantificador
(\cref{prop:b1:logic:negquant}) e usando $\lnot(P \implies Q) \iff
P \land \lnot Q$:
\begin{enumerate}
  \item $\exists x \in \R,\ \forall y \in \R,\ x + y \leq 0$;
  \item $\forall x \in \R,\ \exists y \in \R,\ xy \neq 0$;
  \item $\exists \varepsilon > 0,\ \forall \delta > 0,\ \exists x \in
        \R,\ \abs{x} \leq \delta \text{ e } \abs{f(x)} >
        \varepsilon$.
\end{enumerate}
A \omterm{def:b1:logic:statement}{proposição} (1) é verdadeira: dado $x$, tome $y = -x + 1$; então $x + y = 1 >
0$. A \omterm{def:b1:logic:statement}{proposição} (2) é verdadeira: $x = 0$ satisfaz $xy = 0$ para todo $y$.
\end{solution}

\begin{solution}{exo:b1:logic:2}
Tabela-verdade, escrevendo V/F para os quatro casos $(P, Q)$:
\begin{center}
\begin{tabular}{cc|c|c|c|c}
$P$ & $Q$ & $P \implies Q$ & $\lnot(P \implies Q)$ & $\lnot Q$ &
$P \land \lnot Q$ \\
\hline
V & V & V & F & F & F \\
V & F & F & V & V & V \\
F & V & V & F & F & F \\
F & F & V & F & V & F \\
\end{tabular}
\end{center}
As colunas $4$ e $6$ coincidem, o que demonstra a equivalência. A negação de
``se uma função é derivável, então ela é contínua'' é, portanto:
``existe uma função que é derivável e não é contínua''
(\omterm{def:b1:logic:statement}{proposição} falsa, aliás: a implicação original é verdadeira,
veja o \cref{ch:b1:derivative}).
\end{solution}

\begin{solution}{exo:b1:logic:3}
\emph{Contraposição.} Suponha $x < 1$. Então $x^3 < 1$ (a função cubo
é crescente) e $x < 1$, logo $x^3 + x < 2$. Isso demonstra a
contrapositiva e, portanto, o enunciado.

\emph{Absurdo.} Suponha que $a > 0$ seja o menor real estritamente
positivo. Então $a/2$ é estritamente positivo e $a/2 < a$ (pois $a > 0$),
o que contradiz a minimalidade. Logo, tal $a$ não existe.
\end{solution}

\begin{solution}{exo:b1:logic:4}
\begin{enumerate}
  \item Caso base $n = 0$: $2^0 = 1 = 2^1 - 1$. Passo: supondo a
        identidade para $n$,
        \[
        \sum_{k=0}^{n+1} 2^k = (2^{n+1} - 1) + 2^{n+1}
        = 2 \cdot 2^{n+1} - 1 = 2^{n+2} - 1 .
        \]
  \item Caso base $n = 0$: $4^0 + 5 = 6 = 3 \times 2$. Passo: se
        $4^n + 5 = 3m$, então
        \[
        4^{n+1} + 5 = 4(4^n + 5) - 15 = 3(4m - 5),
        \]
        divisível por $3$.
\end{enumerate}
\end{solution}

\begin{solution}{exo:b1:logic:5}
O passo de indução supõe silenciosamente que os dois grupos (``os $n$
primeiros'' e ``os $n$ últimos'') se sobrepõem, de modo que os lápis
comuns transportem a cor de um grupo ao outro. Para $n + 1 = 2$ os dois
grupos são $\{$primeiro lápis$\}$ e $\{$segundo lápis$\}$: são disjuntos,
e o argumento se rompe. Assim, $P(1) \implies P(2)$ nunca foi demonstrada, e
a indução desmorona --- ainda que $P(n) \implies P(n+1)$ seja
válida para todo $n \geq 2$.
\end{solution}

\begin{solution}{exo:b1:logic:6}
\begin{enumerate}
  \item $x \in A \setminus B \iff x \in A \land x \notin B \iff x \in
        A \land x \in \overline{B} \iff x \in A \cap \overline{B}$.
  \item Usando (1) e a distributividade
        (\cref{prop:b1:logic:setalgebra}):
        $(A \cup B) \cap \overline{C} = (A \cap \overline{C}) \cup
        (B \cap \overline{C})$.
  \item Suponha $A \subseteq B$. Então $A \cup B \subseteq B$ (as duas
        peças estão em $B$) e $B \subseteq A \cup B$ sempre, logo
        $A \cup B = B$. Suponha $A \cup B = B$: então $A \cap B
        \subseteq A$ sempre, e $A \subseteq A \cup B = B$ dá
        $A \subseteq A \cap B$, logo $A \cap B = A$. Suponha $A \cap B =
        A$: então $A = A \cap B \subseteq B$. As três condições são,
        portanto, equivalentes (demonstramos um ciclo de implicações).
\end{enumerate}
\end{solution}

\begin{solution}{exo:b1:logic:7}
\begin{enumerate}
  \item $f$ é \omterm{def:b1:logic:inj}{injetiva} ($n + 1 = m + 1 \implies n = m$), mas não é
        \omterm{def:b1:logic:inj}{sobrejetiva}: $0$ não tem \omterm{def:b1:logic:map}{pré-imagem} em $\N$.
  \item $g$ é \omterm{def:b1:logic:inj}{bijetiva}: $n \mapsto n - 1$ é uma inversa bilateral em
        $\Z$.
  \item $h$ é \omterm{def:b1:logic:inj}{injetiva}: $\frac{x+1}{x-1} = \frac{x'+1}{x'-1}$ dá
        $(x+1)(x'-1) = (x'+1)(x-1)$, isto é, $xx' - x + x' - 1 = xx' -
        x' + x - 1$, logo $2x' = 2x$. Ela não é \omterm{def:b1:logic:inj}{sobrejetiva} sobre $\R$:
        resolver $y = \frac{x+1}{x-1}$ dá $x(y - 1) = y + 1$, que
        não tem solução quando $y = 1$ (a equação fica $0 = 2$).
        Com contradomínio $\R \setminus \{1\}$, o mesmo cálculo dá
        a única \omterm{def:b1:logic:map}{pré-imagem} $x = \frac{y+1}{y-1}$, de modo que
        $h \colon \R \setminus \{1\} \to \R \setminus \{1\}$ é
        \omterm{def:b1:logic:inj}{bijetiva} e $h^{-1}(y) = \frac{y+1}{y-1} = h(y)$: $h$ é a sua
        própria inversa.
\end{enumerate}
\end{solution}

\begin{solution}{exo:b1:logic:8}
\begin{enumerate}
  \item $x \in f^{-1}(B \cap B') \iff f(x) \in B \cap B' \iff f(x) \in
        B \land f(x) \in B' \iff x \in f^{-1}(B) \cap f^{-1}(B')$.
        Para as imagens: $y \in f(A \cup A')$ se, e somente se, $y = f(x)$ para algum $x$
        em $A$ ou em $A'$, ou seja, $y \in f(A)$ ou $y \in f(A')$.
  \item Se $y \in f(A \cap A')$, então $y = f(x)$ com $x \in A$ e
        $x \in A'$, logo $y \in f(A)$ e $y \in f(A')$. Estrita inclusão:
        tome $f \colon \R \to \R$, $x \mapsto x^2$, $A = \{-1\}$,
        $A' = \{1\}$: então $f(A \cap A') = f(\emptyset) = \emptyset$,
        mas $f(A) \cap f(A') = \{1\}$.
  \item ($\Leftarrow$) Com $A = \{x\}$, $A' = \{x'\}$ para $x \neq
        x'$: se $f(x) = f(x')$, então $f(A) \cap f(A') = \{f(x)\}$,
        ao passo que $f(A \cap A') = \emptyset$, o que contradiz a
        igualdade suposta; logo $f$ é \omterm{def:b1:logic:inj}{injetiva}. ($\Rightarrow$) Seja $f$
        \omterm{def:b1:logic:inj}{injetiva} e $y \in f(A) \cap f(A')$: $y = f(x) = f(x')$ com
        $x \in A$, $x' \in A'$; a injetividade dá $x = x' \in A \cap
        A'$, logo $y \in f(A \cap A')$. Junto com (2), vale a igualdade.
\end{enumerate}
\end{solution}

\begin{solution}{exo:b1:logic:9}
$g \circ f = \mathrm{id}_E$ é \omterm{def:b1:logic:inj}{injetiva} e \omterm{def:b1:logic:inj}{sobrejetiva}, logo, pela
\cref{prop:b1:logic:comp}~(2), $f$ é \omterm{def:b1:logic:inj}{injetiva} e $g$ é \omterm{def:b1:logic:inj}{sobrejetiva}.
Exemplo: $E = \N$, $F = \Z$, $f$ a inclusão $n \mapsto n$, e
$g \colon \Z \to \N$, $g(n) = n$ para $n \geq 0$ e $g(n) = 0$ para
$n < 0$. Então $g(f(n)) = n$ para todo $n \in \N$, mas $f$ não é
\omterm{def:b1:logic:inj}{sobrejetiva} e $g$ não é \omterm{def:b1:logic:inj}{injetiva}.
\end{solution}

\begin{solution}{exo:b1:logic:10}
$x^2 - y^2 = x - y \iff (x - y)(x + y) = x - y \iff (x - y)(x + y - 1)
= 0 \iff y = x$ ou $y = 1 - x$. \emph{Reflexiva:} $y = x$ serve.
\emph{Simétrica:} a condição ``$y = x$ ou $y = 1 - x$'' é
simétrica em $x$ e $y$ (se $y = 1 - x$, então $x = 1 - y$).
\emph{Transitiva:} suponha $x \mathbin{\mathcal{R}} y$ e $y
\mathbin{\mathcal{R}} z$; percorrendo os quatro casos, $z$ é igual a $x$
ou a $1 - x$ em cada um deles (por exemplo, $y = 1 - x$ e $z = 1 - y$ dão $z = x$).
Logo $\mathcal{R}$ é uma \omterm{def:b1:logic:equiv}{relação de equivalência} e $\mathrm{cl}(x) =
\{x,\, 1 - x\}$. Essa classe tem um só elemento exatamente quando $x = 1 - x$,
isto é, para $x = \frac12$.
\end{solution}

\begin{solution}{exo:b1:logic:11}
Seja $f \colon E \to \mathcal{P}(E)$ uma \omterm{def:b1:logic:map}{aplicação} qualquer e ponha
$D = \{x \in E : x \notin f(x)\} \in \mathcal{P}(E)$. Suponha $D =
f(a)$ para algum $a \in E$. Se $a \in D$, então, por definição de $D$,
$a \notin f(a) = D$: contradição. Se $a \notin D$, então $a \notin
f(a)$, e por definição de $D$, $a \in D$: contradição. Logo $D$ não
está na imagem de $f$, e $f$ não é \omterm{def:b1:logic:inj}{sobrejetiva}. (Em particular, nenhum
\omterm{def:b1:logic:sets}{conjunto} está em bijeção com o seu \omterm{def:b1:logic:sets}{conjunto das partes}: há ``mais''
subconjuntos de $\N$ do que inteiros.)
\end{solution}

\begin{solution}{exo:b1:logic:12}
\begin{enumerate}
  \item ($\Rightarrow$) Seja $f$ \omterm{def:b1:logic:inj}{sobrejetiva} e $\Phi(B) =
        \Phi(B')$. Para $y \in B$, escolha $x$ com $f(x) = y$; então $x
        \in f^{-1}(B) = f^{-1}(B')$, logo $y = f(x) \in B'$. Portanto $B
        \subseteq B'$, e simetricamente $B' \subseteq B$: $\Phi$ é
        \omterm{def:b1:logic:inj}{injetiva}. ($\Leftarrow$) Se $f$ não é \omterm{def:b1:logic:inj}{sobrejetiva}, escolha $y_0
        \in F$ fora da imagem; então $f^{-1}(\{y_0\}) = \emptyset =
        f^{-1}(\emptyset)$ com $\{y_0\} \neq \emptyset$, de modo que $\Phi$ não
        é \omterm{def:b1:logic:inj}{injetiva}.
  \item ($\Rightarrow$) Seja $f$ \omterm{def:b1:logic:inj}{injetiva} e $A \subseteq E$. Ponha
        $B = f(A)$; então $f^{-1}(B) = \{x : f(x) \in f(A)\}$, e a
        injetividade dá $f(x) \in f(A) \iff x \in A$, logo $\Phi(B) =
        A$: $\Phi$ é \omterm{def:b1:logic:inj}{sobrejetiva}. ($\Leftarrow$) Se $f$ não é
        \omterm{def:b1:logic:inj}{injetiva}, tome $x \neq x'$ com $f(x) = f(x')$. Toda
        \omterm{def:b1:logic:map}{pré-imagem} $f^{-1}(B)$ contém $x$ se, e somente se, contém
        $x'$; logo $\{x\}$ não é da forma $\Phi(B)$, e
        $\Phi$ não é \omterm{def:b1:logic:inj}{sobrejetiva}.
\end{enumerate}
\end{solution}

\begin{solution}{pb:b1:logic:1}
\textbf{1.} \emph{Reflexiva:} $\mathrm{id}_E$ é uma bijeção de $E$
sobre si mesmo. \emph{Simétrica:} se $f \colon E \to F$ é \omterm{def:b1:logic:inj}{bijetiva}, o
\cref{thm:b1:logic:inverse} fornece $f^{-1} \colon F \to E$, ela própria
\omterm{def:b1:logic:inj}{bijetiva}. \emph{Transitiva:} se $f \colon E \to F$ e $g \colon F
\to G$ são bijeções, a \cref{prop:b1:logic:comp}~(1) diz que $g \circ f
\colon E \to G$ é uma bijeção. (Isso é apenas ``como'' uma \omterm{def:b1:logic:equiv}{relação de
equivalência}: a coleção de todos os \omterm{def:b1:logic:sets}{conjuntos} não é ela própria um
\omterm{def:b1:logic:sets}{conjunto}, pelos paradoxos que o \cref{exo:b1:logic:11} insinua; o que importa
são as três propriedades.)

\textbf{2.} Se $f \colon E \to F$ e $g \colon F \to G$ são
\omterm{def:b1:logic:inj}{injetivas}, $g \circ f$ é \omterm{def:b1:logic:inj}{injetiva} pela \cref{prop:b1:logic:comp}~(1):
$E \preceq G$. Para o segundo ponto, correstrinja $f$ à sua imagem: a
\omterm{def:b1:logic:map}{aplicação} $\tilde f \colon E \to f(E)$, $x \mapsto f(x)$, é \omterm{def:b1:logic:inj}{sobrejetiva} por
construção de $f(E)$ e \omterm{def:b1:logic:inj}{injetiva} porque $f$ o é, logo \omterm{def:b1:logic:inj}{bijetiva}:
$E \approx f(E)$.

\textbf{3.} ($\Rightarrow$) Seja $f \colon E \to F$ \omterm{def:b1:logic:inj}{injetiva} e
fixe $a \in E$ ($E \neq \emptyset$). Defina $s \colon F \to E$ por:
$s(y)$ é o único $x$ com $f(x) = y$ quando $y \in f(E)$
(unicidade pela injetividade), e $s(y) = a$ caso contrário. Para todo $x
\in E$, $s(f(x)) = x$, logo todo $x$ é atingido: $s$ é \omterm{def:b1:logic:inj}{sobrejetiva}.
($\Leftarrow$) Seja $s \colon F \to E$ \omterm{def:b1:logic:inj}{sobrejetiva}. Para cada $x \in
E$, escolha um $y_x \in F$ com $s(y_x) = x$ e ponha $u(x) = y_x$. Se
$u(x) = u(x')$, então $x = s(u(x)) = s(u(x')) = x'$: $u \colon E \to F$
é \omterm{def:b1:logic:inj}{injetiva}.

\textbf{4.} $n \mapsto n + 1$ leva $\N$ em $\N^*$, é \omterm{def:b1:logic:inj}{injetiva}
($n + 1 = m + 1 \implies n = m$) e \omterm{def:b1:logic:inj}{sobrejetiva} (todo $m \geq 1$ é
$(m - 1) + 1$ com $m - 1 \in \N$). Quanto a $\sigma$: ele leva os números
pares $0, 2, 4, \dots$ em $0, 1, 2, \dots$ e os ímpares $1, 3,
5, \dots$ em $-1, -2, -3, \dots$ Injetividade: as entradas pares caem em
$\N$ ($\sigma(n) = n/2 \geq 0$) e as entradas ímpares caem nos inteiros
estritamente negativos ($\sigma(n) = -(n+1)/2 \leq -1$), de modo que uma
colisão teria de ocorrer dentro de uma mesma classe de paridade, onde
$\sigma$ é estritamente monótona ($n/2 = m/2$ ou $(n+1)/2 = (m+1)/2$ força $n =
m$). Sobrejetividade: $k \geq 0$ é $\sigma(2k)$; $k \leq -1$ é
$\sigma(-2k - 1)$ com $-2k - 1 \geq 1$ ímpar. Logo $\N \approx \N^*$ e
$\N \approx \Z$.

\textbf{5.} $C_0 = E \setminus g(F) \subseteq C$, de modo que $x \notin C$
implica $x \notin C_0$, isto é, $x \in g(F)$: algum $y \in F$ satisfaz
$g(y) = x$. Se também $g(y') = x$, a injetividade de $g$ dá $y' = y$.
Assim, a segunda cláusula da definição de $h$ seleciona um único
elemento, bem definido, $g^{-1}(x)$.

\textbf{6.} As imagens diretas comutam com as uniões
(\cref{exo:b1:logic:8}~(1), aplicado a $f$ e depois a $g$):
\[
g\bigl(f(C)\bigr)
= g\Bigl(f\Bigl(\bigcup_{n \in \N} C_n\Bigr)\Bigr)
= \bigcup_{n \in \N} g\bigl(f(C_n)\bigr)
= \bigcup_{n \in \N} C_{n+1}
= \bigcup_{n \geq 1} C_n \subseteq C .
\]

\textbf{7.} Sejam $x \neq x'$ em $E$. Se ambos estão em $C$, então $h(x) =
f(x) \neq f(x') = h(x')$ pela injetividade de $f$. Se nenhum dos dois está em
$C$, então $g(h(x)) = x \neq x' = g(h(x'))$, logo $h(x) \neq h(x')$. Se
$x \in C$ e $x' \notin C$ (o caso misto, a menos de troca de nomes):
suponha $h(x) = h(x')$, isto é, $f(x) = g^{-1}(x')$. Aplicando $g$:
$x' = g(f(x)) \in g(f(C))$, e a questão 6 dá $x' \in C$ ---
contradição. Logo $h(x) \neq h(x')$ em todos os casos: $h$ é \omterm{def:b1:logic:inj}{injetiva}.

\textbf{8.} Seja $y \in F$. \emph{Caso 1: $g(y) \notin C$.} Então
$h(g(y)) = g^{-1}(g(y)) = y$: o elemento $g(y)$ é uma \omterm{def:b1:logic:map}{pré-imagem}.
\emph{Caso 2: $g(y) \in C$}, digamos $g(y) \in C_n$. Como $g(y) \in
g(F)$, temos $g(y) \notin C_0 = E \setminus g(F)$, logo $n \geq 1$
e $g(y) \in C_n = g(f(C_{n-1}))$: existe $x \in C_{n-1}$ com
$g(y) = g(f(x))$. A injetividade de $g$ dá $y = f(x)$, e $x \in
C_{n-1} \subseteq C$, logo $h(x) = f(x) = y$. Nos dois casos, $y$ é
atingido: $h$ é \omterm{def:b1:logic:inj}{sobrejetiva} e, portanto, \omterm{def:b1:logic:inj}{bijetiva}.

\textbf{9.} Se $E \preceq F$ e $F \preceq E$, escolha injeções $f
\colon E \to F$ e $g \colon F \to E$; as questões 5--8 constroem uma
bijeção $h \colon E \to F$, logo $E \approx F$. O enunciado é
não trivial porque as duas injeções dadas não guardam relação alguma ---
nenhuma delas precisa ser \omterm{def:b1:logic:inj}{sobrejetiva}, e nenhuma fórmula ingênua que
misture $f$ e $g$ define uma \omterm{def:b1:logic:map}{aplicação}: todo o conteúdo está na \omterm{thm:b1:logic:partition}{partição}
de $E$ na região $C$ (onde se copia $f$) e no seu complementar (onde se
percorre $g$ ao contrário).

\textbf{10.} (a) A inclusão $\intoo01 \to \intcc01$ é \omterm{def:b1:logic:inj}{injetiva};
e $x \mapsto \frac{x + 1}3$ leva $\intcc01$ injetivamente em
$\intcc{\frac13}{\frac23} \subseteq \intoo01$ (é afim com
coeficiente angular não nulo). Pela questão 9, $\intcc01 \approx \intoo01$ --- uma
bijeção bastante desagradável de escrever explicitamente.
(b) \emph{Injetividade.} Suponha $2^p(2q + 1) = 2^{p'}(2q' + 1)$ com,
digamos, $p \leq p'$. Dividindo por $2^p$: $2q + 1 = 2^{p' - p}(2q' + 1)$.
Se $p' > p$, o lado direito é par e o esquerdo é ímpar ---
impossível; logo $p = p'$, e então $2q + 1 = 2q' + 1$ e $q = q'$.
\emph{Sobrejetividade.} Mostramos por indução forte que todo inteiro
$m \geq 1$ é da forma $2^p(2q + 1)$. Para $m = 1$: $p = q = 0$.
Seja $m \geq 1$ e suponha a afirmação para todos os inteiros de
$\intint1m$. Se $m + 1$ é ímpar, $m + 1 = 2q + 1$ com $p = 0$. Se
$m + 1$ é par, $m + 1 = 2m'$ com $1 \leq m' \leq m$; por
hipótese, $m' = 2^p(2q + 1)$, logo $m + 1 = 2^{p+1}(2q + 1)$. Portanto
$\varphi(p, q) = 2^p(2q + 1) - 1$ atinge todo $n \in \N$, e
$\varphi$ é uma bijeção $\N \times \N \to \N$.

\textbf{11.} Como $A$ é infinito, $A \setminus \{\varphi(0), \dots,
\varphi(n - 1)\}$ nunca é vazio, e a propriedade do menor elemento de
$\N$ (usada para demonstrar o \cref{thm:b1:logic:induction}) torna legítima a
definição recursiva. \emph{Estritamente crescente:}
$\varphi(n + 1)$ pertence a $A \setminus \{\varphi(0), \dots,
\varphi(n)\} \subseteq A \setminus \{\varphi(0), \dots, \varphi(n -
1)\}$, cujo mínimo é $\varphi(n)$; logo $\varphi(n + 1) \geq
\varphi(n)$, e a igualdade está excluída, donde $\varphi(n+1) >
\varphi(n)$. \emph{$\varphi(n) \geq n$:} por indução, $\varphi(0)
\geq 0$, e $\varphi(n + 1) \geq \varphi(n) + 1 \geq n + 1$.
A \emph{injetividade} decorre da monotonicidade estrita.
\emph{Sobrejetividade sobre $A$:} suponha que algum $a \in A$ nunca seja
atingido. Como $\varphi(a + 1) \geq a + 1 > a$, o \omterm{def:b1:logic:sets}{conjunto} dos $n$ com
$\varphi(n) > a$ é não vazio; seja $n$ o seu menor elemento. Para todo
$k < n$, $\varphi(k) \leq a$, e portanto $\varphi(k) < a$ ($a$ não é
atingido). Então $a$ está em $A \setminus \{\varphi(0), \dots,
\varphi(n - 1)\}$ e $a < \varphi(n)$, contradizendo a minimalidade
que define $\varphi(n)$. Logo $\varphi$ é uma bijeção $\N \to A$, e
$A$ é enumerável.

\textbf{12.} Seja $E \preceq \N$ por meio de uma injeção $f$; então $E \approx
f(E)$ (questão 2). Se $f(E)$ é finito, $E$ é finito; se $f(E)$ é
infinito, a questão 11 dá $f(E) \approx \N$, logo $E \approx \N$ por
transitividade (questão 1). Reciprocamente, \omterm{def:b1:logic:sets}{conjuntos} finitos e \omterm{pb:b1:logic:1}{conjuntos
enumeráveis} injetam-se obviamente em $\N$. O atalho: $E \preceq \N$ e $\N
\preceq E$ dão $E \approx \N$ diretamente por
Cantor--Schröder--Bernstein --- sem nenhum argumento de enumeração.

\textbf{13.} Sejam $f \colon E \to \N$ e $g \colon F \to \N$
injeções. Então $(x, y) \mapsto \varphi\bigl(f(x), g(y)\bigr)$ é uma
injeção $E \times F \to \N$: se as imagens coincidem, a injetividade de
$\varphi$ (questão 10) dá $f(x) = f(x')$ e $g(y) = g(y')$,
e então $x = x'$, $y = y'$. Quanto a $\Z \times \N^*$: os dois fatores são
enumeráveis (questão 4), logo $\Z \times \N^* \preceq \N$; ele é
infinito (contém $\{0\} \times \N^*$) e, portanto, enumerável pela
questão 12.

\textbf{14.} Todo racional $r$ tem uma única representação $r =
p/q$ com $p \in \Z$, $q \in \N^*$ e a fração irredutível
(a unicidade é demonstrada no \cref{ch:b1:arith}; para $r = 0$, tome
$0/1$). A \omterm{def:b1:logic:map}{aplicação} $r \mapsto (p, q)$ é então \omterm{def:b1:logic:inj}{injetiva}: o par
determina $r = p/q$. Portanto $\Q \preceq \Z \times \N^* \preceq \N$
pela questão 13. Como $\N \subseteq \Q$ dá $\N \preceq \Q$,
a questão 12 (ou diretamente Cantor--Schröder--Bernstein) mostra que $\Q
\approx \N$: os racionais são enumeráveis.

\textbf{15.} Para cada $n$, fixe uma injeção $f_n \colon E_n \to \N$.
Para $x \in \bigcup_n E_n$, seja $n(x)$ o \emph{menor} $n$ com
$x \in E_n$ e ponha $u(x) = \varphi\bigl(n(x), f_{n(x)}(x)\bigr)
\in \N$. Se $u(x) = u(x')$, a injetividade de $\varphi$ dá $n(x) =
n(x') = n$ e $f_n(x) = f_n(x')$, e então $x = x'$ pela injetividade de
$f_n$. Assim, a união se injeta em $\N$: ela é no máximo enumerável.

\textbf{16.} Seja $\Psi(F) = \sum_{i \in F} 2^i$ para $F \subseteq \N$
finito ($\Psi(\emptyset) = 0$). Suponha $F \neq F'$ e seja $m$
o maior elemento em que eles diferem, digamos $m \in F \setminus
F'$ (troque
os nomes se necessário). Os elementos $> m$ pertencem a ambos ou a
nenhum, de modo que contribuem igualmente para as duas somas; comparando
as contribuições dos elementos $\leq m$:
\[
\sum_{i \in F,\, i \leq m} 2^i \geq 2^m
> 2^m - 1 = \sum_{k=0}^{m-1} 2^k
\geq \sum_{i \in F',\, i \leq m} 2^i ,
\]
usando a soma geométrica do \cref{exo:b1:logic:4}. Portanto $\Psi(F)
\neq \Psi(F')$: $\Psi$ é \omterm{def:b1:logic:inj}{injetiva} e o \omterm{def:b1:logic:sets}{conjunto} dos subconjuntos finitos
de $\N$ é no máximo enumerável; ele é infinito (contém todos os
\omterm{def:b1:logic:sets}{conjuntos} unitários) e, portanto, enumerável.

\textbf{17.} Envie $A \subseteq \N$ à sua função indicadora $\mathbf 1_A
\colon \N \to \{0,1\}$, $\mathbf 1_A(n) = 1$ se $n \in A$ e $0$
caso contrário; envie $u \in \{0,1\}^{\N}$ a $A_u = \{n \in \N : u(n) =
1\}$. As duas aplicações são inversas uma da outra: $A_{\mathbf 1_A} = A$ e
$\mathbf 1_{A_u} = u$ (verifique o valor em cada $n$). Pelo
\cref{thm:b1:logic:inverse}, cada uma é uma bijeção: $\mathcal{P}(\N)
\approx \{0,1\}^{\N}$.

\textbf{18.} Para todo $n$, $d(n) = 1 - \Phi(n)(n) \neq \Phi(n)(n)$,
de modo que as sequências $d$ e $\Phi(n)$ diferem no índice $n$: $d \neq
\Phi(n)$. Logo nenhum $\Phi$ é \omterm{def:b1:logic:inj}{sobrejetivo} e, pela questão 3, também não
existe injeção $\{0,1\}^{\N} \to \N$: $\{0,1\}^{\N}$ não é no
máximo enumerável. Através do dicionário da questão 17, uma \omterm{def:b1:logic:map}{aplicação}
$\Phi
\colon \N \to \{0,1\}^{\N}$ é uma \omterm{def:b1:logic:map}{aplicação} $f \colon \N \to
\mathcal{P}(\N)$, e $d$ corresponde ao \omterm{def:b1:logic:sets}{conjunto} $D = \{n : n
\notin f(n)\}$ (com efeito, $d(n) = 1 \iff \Phi(n)(n) = 0 \iff n \notin
f(n)$): o argumento diagonal \emph{é} a demonstração de Cantor do
\cref{exo:b1:logic:11} para $E = \N$.

\textbf{19.} Escreva $x_n = 0.d_1(n)\,d_2(n)\,d_3(n)\dots$ na forma
própria e defina $\delta_n = 5$ se $d_n(n) \neq 5$, $\delta_n = 6$
se $d_n(n) = 5$, e depois $x = 0.\delta_1\delta_2\delta_3\dots$ Essa
expansão usa apenas os algarismos $5$ e $6$, de modo que não termina em
uma cadeia de $9$s: é a expansão própria de um real $x \in \intco01$.
Para cada $n$, os $n$-ésimos algarismos de $x$ e de $x_n$ diferem
($\delta_n \neq d_n(n)$, por construção); como as expansões próprias
são únicas, $x \neq x_n$. Assim, nenhuma sequência esgota $\intco01$: pela
questão 3, novamente, $\intco01$ não é no máximo enumerável.

\textbf{20.} $\intco01 \subseteq \R$, de modo que uma injeção $\R \to \N$
se restringiria a uma injeção em $\intco01$, contradizendo a questão 19:
$\R$ é não enumerável. Se $\R \setminus \Q$ fosse no máximo enumerável,
então $\R = \Q \cup (\R \setminus \Q)$ seria uma união de dois
\omterm{def:b1:logic:sets}{conjuntos} no máximo enumeráveis e, portanto, no máximo enumerável pela
questão 15 (tome $E_0 = \Q$, $E_n = \R \setminus \Q$ para $n \geq 1$) ---
contradição. Logo os irracionais são não enumeráveis. Mais precisamente:
dentro de $\R$, os racionais formam um \omterm{pb:b1:logic:1}{conjunto enumerável}, enquanto o
seu complementar é não enumerável; nenhuma bijeção pode jamais casar
$\R
\setminus \Q$ com $\Q$ --- há estritamente ``mais'' irracionais
do que racionais, embora ambos sejam infinitos e ambos sejam densos.

\textbf{21.} $p/q$ (com $q \neq 0$) é raiz de $qX - p$, um
polinômio não nulo com coeficientes inteiros. $\sqrt 2$ é raiz
de $X^2 - 2$. Para $x = \sqrt 2 + \sqrt 3$: $x^2 = 5 + 2\sqrt 6$, logo
$x^2 - 5 = 2\sqrt 6$ e $(x^2 - 5)^2 = 24$, isto é,
\[
x^4 - 10x^2 + 1 = 0 :
\]
$\sqrt 2 + \sqrt 3$ é raiz de $X^4 - 10X^2 + 1$.

\textbf{22.} Associe a $P = a_0 + a_1X + \dots + a_nX^n$ (grau $\leq n$,
coeficientes inteiros) a lista $(a_0, \dots, a_n) \in \Z^{n+1}$: isso é
\omterm{def:b1:logic:inj}{injetivo}, pois um polinômio fica determinado por seus coeficientes. Por
indução em $n$: $\Z^1 = \Z$ é enumerável (questão 4), e
$\Z^{n+2} \approx \Z^{n+1} \times \Z$ é no máximo enumerável pela
questão 13. Logo cada \omterm{def:b1:logic:sets}{conjunto} de polinômios inteiros de grau limitado é
no máximo enumerável; ele é infinito (contém as constantes) e, portanto,
enumerável pela questão 12.

\textbf{23.} O \omterm{def:b1:logic:sets}{conjunto} de todos os polinômios inteiros é $\bigcup_{n \in
\N} \{P : \deg P \leq n,\ P \text{ has integer coefficients}\}$, uma
união enumerável de \omterm{pb:b1:logic:1}{conjuntos enumeráveis}: no máximo enumerável pela
questão 15, infinito e, portanto, enumerável.

\textbf{24.} Para cada polinômio inteiro não nulo $P$, o \omterm{def:b1:logic:sets}{conjunto} de raízes
$R_P = \{x \in \R : P(x) = 0\}$ é finito (no máximo $\deg P$
elementos, admitido). Pela questão 23, os polinômios inteiros não nulos
podem ser enumerados $P_0, P_1, P_2, \dots$; então $\mathcal{A} =
\bigcup_{n \in \N} R_{P_n}$ é uma união enumerável de \omterm{def:b1:logic:sets}{conjuntos} finitos (logo
no máximo enumeráveis): no máximo enumerável pela questão 15. Ele
contém $\Q$ (questão 21), logo é infinito: $\mathcal{A}$ é
enumerável.

\textbf{25.} Se $\R \setminus \mathcal{A}$ fosse no máximo enumerável,
$\R = \mathcal{A} \cup (\R \setminus \mathcal{A})$ seria no máximo
enumerável (questão 15), contradizendo a questão 20. Logo existem
\omterm{pb:b1:logic:1}{números transcendentes}, que formam até mesmo um \omterm{def:b1:logic:sets}{conjunto} não enumerável,
ao passo que os \omterm{pb:b1:logic:1}{números algébricos} --- entre os quais figura todo número
construído a partir de inteiros por radicais --- formam um mero esqueleto
enumerável dentro de $\R$. Resumo da arquitetura: as questões 1--3 montam
a linguagem da comparação; Cantor--Schröder--Bernstein (questões
5--9) permite demonstrar a equipotência por meio de duas injeções fáceis,
em vez de uma bijeção engenhosa, e foi usado para $\intcc01 \approx \intoo01$,
para $\Q$ e ao longo de toda a Parte V; a bijeção de emparelhamento
(questão 10) alimentou os produtos e as uniões enumeráveis (questões 13
e 15), que por sua vez alimentaram $\Q$, os polinômios inteiros e
$\mathcal{A}$; o argumento diagonal (questões 18--19) forneceu
a única desigualdade estrita $\N \prec \R$ que torna toda a história
não trivial. A conclusão de Cantor é filosoficamente notável: a
demonstração não exibe nenhum \omterm{pb:b1:logic:1}{número transcendente}, e ainda assim mostra
que, no sentido da equipotência, \emph{quase todo} número real é
transcendente. Nomear um transcendente específico --- $\pi$ ou
$\eu$ --- exigiu uma matemática inteiramente diferente e décadas a mais
de trabalho.
\end{solution}
